\usepackage{xcolor}
\definecolor{ugblue}{RGB}{31, 78, 121}
\usepackage{graphicx}
\usepackage{mathptmx} % Times New Roman
\usepackage[T1]{fontenc}
\usepackage{setspace}
\setstretch{1.5} % Espaciado 1.5 para el cuerpo del trabajo
\setlength{\parskip}{8pt}
\setlength{\parindent}{0pt}

% Configuración de comillas APA (estándar en lugar de angulares)
\usepackage[spanish=mexican]{csquotes}
\DeclareQuoteAlias{mexican}{spanish}
\AtBeginDocument{\setquotestyle{mexican}}

% Numeración inferior derecha
\usepackage{fancyhdr}
\pagestyle{fancy}
\fancyhf{}
\renewcommand{\headrulewidth}{0pt}
\fancyfoot[R]{\thepage}

% Ajuste de tamaños para títulos y subtítulos (14pt)
\usepackage{titlesec}
\titleformat{\section}{\fontsize{14}{17}\bfseries}{\thesection}{1em}{}
\titleformat{\subsection}{\fontsize{14}{17}\bfseries}{\thesubsection}{1em}{}
\titleformat{\subsubsection}{\fontsize{14}{17}\bfseries}{\thesubsubsection}{1em}{}

% Caption APA 7: normalsize, doble espaciado, número en negrita, título en cursiva en nueva línea
\usepackage{caption}
\usepackage{floatrow}
\floatsetup[figure]{capposition=top}
\floatsetup[table]{capposition=top}
\captionsetup{
  labelsep        = newline,
  font            = {normalsize, doublespacing},
  labelfont       = bf,
  textfont        = it,
  justification   = raggedright,
  singlelinecheck = false,
  skip            = 6pt
}

% Asegurar interlineado doble en las notas de threeparttable (kableExtra)
\usepackage{threeparttable}
\usepackage{threeparttablex}
\usepackage{etoolbox}
% Reserva de espacio vertical. Una tabla fijada con [H] y su Nota son dos
% bloques distintos para LaTeX, de modo que un salto de pagina puede caer
% entre ambos y dejar la nota huerfana. \Needspace obliga a romper antes
% de la tabla cuando no cabe el bloque entero.
\usepackage{needspace}
\appto\TPTnoteSettings{\linespread{2}\selectfont}

% Redefinir el índice como una sección numerada llamada "Tabla de contenido"
\AtBeginDocument{%
  \renewcommand{\contentsname}{Tabla de contenido}%
  \renewcommand{\tablename}{Tabla}%   % sobrescribe la traducción de babel
}
\makeatletter
\renewcommand\tableofcontents{%
    \addtocontents{toc}{\protect\setcounter{tocdepth}{-1}}%
    \section{\contentsname}%
    \addtocontents{toc}{\protect\setcounter{tocdepth}{3}}%
    \@starttoc{toc}%
}
\makeatother

% ── Tablas Likert del instrumento (Anexo II) ─────────────────────────────────
% El Anexo II es un facsímil del cuestionario, no una tabla de resultados: la
% afirmación tiene que caber en una o dos líneas y las casillas 1-5 tienen que
% poder marcarse. Antes iban como tabla Markdown, que reparte el ancho a partes
% iguales entre las siete columnas y dejaba la afirmación tan estrecha que cada
% ítem se partía en cuatro a seis líneas. Se fija aquí y no en cada tabla para
% que las tres salgan idénticas.
%
% Tres decisiones y su porqué:
%   - `spacing{1}`: el cuerpo va a 1,5 y ese interlineado se heredaba dentro de
%     las celdas, multiplicando por 1,5 la altura de cada línea partida.
%   - `footnotesize`: son 10 pt sobre un cuerpo de 12, el mismo tamaño que
%     `font_size = 10` da al resto de tablas del documento.
%   - `\square`: es el glifo con el que pandoc compone las casillas `- [ ]` del
%     resto del instrumento, de modo que las dos formas de marcar coinciden.
\newcommand{\lkbox}{$\square$}
\newcommand{\lkfila}[2]{#1 & #2 & \lkbox & \lkbox & \lkbox & \lkbox & \lkbox \\}
\newenvironment{likert}{%
  \par\addvspace{4pt}%
  \begin{spacing}{1}%
  \footnotesize
  \renewcommand{\arraystretch}{1.3}%
  \noindent
  \tabular{@{}l>{\raggedright\arraybackslash}p{9.5cm}ccccc@{}}
  \toprule
  \textbf{N.\textsuperscript{o}} & \textbf{Afirmación} &
  \textbf{1} & \textbf{2} & \textbf{3} & \textbf{4} & \textbf{5} \\
  \midrule
}{%
  \bottomrule
  \endtabular
  \end{spacing}%
  \par\addvspace{12pt}%
}

% ── Anexos en PDF ────────────────────────────────────────────────────────────
% Los anexos administrativos (autorización institucional, aprobación del comité,
% carta del director, informe de similitud) son PDF firmados y escaneados que se
% insertan tal cual. `pdfpages` los intercala; sin él, `\includepdf` no existe.
%
% Las páginas insertadas NO heredan el estilo de página, de modo que perderían la
% numeración inferior derecha que exige el formato de la Universidad. Se recupera
% con `pagecommand={\thispagestyle{fancy}}` en cada llamada; la definición está
% en la función `anexo_pdf()` del chunk `setup`, para no repetirla.
\usepackage{pdfpages}

% Recorta el índice desde el punto donde se invoque: a partir de ahí solo se
% listan las entradas de nivel `section` y `subsection` (cada anexo es una
% `subsection` de «Anexos»; las secciones del cuestionario, más hondas, no).
% Se usa una sola vez, justo antes de la página separadora de los anexos, porque las secciones del instrumento son
% subdivisiones de un documento reproducido y no del trabajo: listarlas hacía
% que el índice pareciera tener siete apartados más.
%
% Va como macro y no escrito en el .Rmd a propósito: pandoc no anida bien los
% dos niveles de llaves de `\addtocontents{toc}{\protect\setcounter{...}{...}}`
% y la compilación aborta con «No counter 'D' defined».
\newcommand{\indicesolosecciones}{%
  \addtocontents{toc}{\protect\setcounter{tocdepth}{2}}%
}
